% RLJ main.tex Version 2025.1

\documentclass[10pt]{article} % For LaTeX2e

%%%%%%%%%%%%%%%%%%%%%%%%%%%%%%%%%%%%%%%%%%%%%%%%%%%%%%%%%%%%%%%%
% AUTHOR: Select ONE option:
%      [accepted]{rlj} --> for camera ready (after peer review, if accepted)
%      {rlj}           --> for submission
%      [preprint]{rlj} --> to de-anonymize and remove references to RLJ/RLC
%%%%%%%%%%%%%%%%%%%%%%%%%%%%%%%%%%%%%%%%%%%%%%%%%%%%%%%%%%%%%%%%
\usepackage{rlj}           % Should be uncommented for submission
%\usepackage[accepted]{rlj} % Should be uncommented for the camera-ready
%\usepackage[preprint]{rlj} % Should be uncommented for preprint versions

%%%%%%%%%%%%%%%%%%%%%%%%%%%%%%%%%%%%%%%%%%%%%%%%%%%%%%%%%%%%%%%%
% WARNING: The following packages are already included in the
%          rlj.sty style file:
%
%  1. fancyhdr  - For controlling header/footers
%  2. natbib    - For formatting the bibliography
%  3. enumitem  - To customize the appearance of lists
%  4. fontenc (with option [T1]) - Allows for proper hyphenation and accents
%  5. times     - Times new roman font
%  6. ragged2e  - Used to justify text
%  7. tcolorbox - Used to create boxes on cover page
%  8. hyperref  - Configures hyperlinks throughout (e.g., links to references)
%  9. xcolor    - Used to define custom colors for links and boxes
%  10. amsmath  - Not used, but conflicts with lineno, so we include (and patch) it for authors
%  11. etoolbox - Included in the amsmath + lineno patch
%  12. lineno   - For adding line numbers when in submission
%
% You do not need to include them again in your main.tex.
% Including them again may lead to conflicts or compilation errors.
% Additionally, avoid loading packages that might conflict with these.
%%%%%%%%%%%%%%%%%%%%%%%%%%%%%%%%%%%%%%%%%%%%%%%%%%%%%%%%%%%%%%%%

%%%%%%%%%%%%%%%%%%%%%%%%%%%%%%%%%%%%%%%%%%%%%%%%%%%%%%%%%%%%%%%%
% Recommended (but not required) packages
%%%%%%%%%%%%%%%%%%%%%%%%%%%%%%%%%%%%%%%%%%%%%%%%%%%%%%%%%%%%%%%%
\usepackage{amssymb}            % Defines common symbols like \mathbb R
\usepackage{mathtools}          % Extends amsmath, providing common math tools
\usepackage{mathrsfs}           % Enables \mathscr, which can work in cases that \mathcal does not
%\mathtoolsset{showonlyrefs}     % Only number equations that are referenced (optional)
\usepackage{graphicx}           % For including images
\usepackage{subcaption}         % Allows for the use of subfigures and subcaptions
\usepackage[space]{grffile}     % For spaces in image names
\usepackage{url}                % For displaying URLs
\usepackage{lipsum}             % For placeholder text

%%%%%%%%%%%%%%%%%%%%%%%%%%%%%%%%%%%%%%%%%%%%%%%%%%%%%%%%%%%%%%%%
% AUTHOR: Fill in the following meta-data
%%%%%%%%%%%%%%%%%%%%%%%%%%%%%%%%%%%%%%%%%%%%%%%%%%%%%%%%%%%%%%%%

% Enter the title of your paper:
\title{Federated Q Learning with Hyperdimensional  Approximation }

% The "running title" will be displayed in the header on every-other page.
% It is typically either the same as the title or a shorter version of the title.
% Enter your running title here:
\setrunningtitle{FedQHD: Federated Q Learning with Hyperdimensional Approximation}

% WARNING: Authors must not appear in the submitted version. They should be hidden
% as long as the rlj package is used without the [accepted] or [preprint] options.
% Non-anonymous submissions will be rejected without review.

% Enter the author names below.
% NOTE: Denote affiliations using superscripts as in the provided example.
% NOTE: Use \textscript{1,2,3} instead of $^{1,2,3}$.
%       - Failure to do so will cause affiliation superscripts to appear on the cover page for camera-ready and preprint versions.
\author{Author Name\textsuperscript{1}}

% NOTE: For camera-ready and preprint versions, the cover page includes author names but not affiliations.
% It automatically removes the superscripts for affiliations.
% If the automatic process breaks (e.g., if an author name should include a superscript), you can manually define the author string to appear on the cover page by uncommenting the following line.
%\coverPageAuthor{Author Name}

% Author emails, which can be clustered if they have shared endings as in this example
\emails{author@institution.edu}

% Author affiliations, in the order the occur
% The inclusion of state/province, etc. is optional.
% The inclusion of multiple affiliations is optional.
%   - List multiple affiliations with comma-separated numbers as in the example.
\affiliations{
$^{1}$\textbf{Institution Name, Department}
}

\contribution{
    % Contribution
    [Describe your first key contribution here]
    }
    {
    % Caveat/Context:
    None
    }

\contribution{
    % Contribution
    [Describe your second key contribution here]
    }
    {
    % Caveat/Context:
    None
    }

% Include a list of keywords for the topic of the paper:
\keywords{Federated Learning, Q-Learning, Hyperdimensional Computing, Reinforcement Learning} % Your keywords

% Define the summary that appears on the cover page.
\summary{This paper investigates federated Q-learning with hyperdimensional approximation (FedQHD) for distributed reinforcement learning. [Add your summary here describing the problem, approach, and key findings.]
}

%%%%%%%%%%%%%%%%%%%%%%%%%%%%%%%%%%%%%%%%%%%%%%%%%%%%%%%%%%%%%%%%
%% Begin document, create title and abstract
%%%%%%%%%%%%%%%%%%%%%%%%%%%%%%%%%%%%%%%%%%%%%%%%%%%%%%%%%%%%%%%%
\begin{document}

\makeCover  % Create the cover page
\maketitle  % Make the title section

\begin{abstract}
[Write your abstract here. The abstract should concisely summarize: (1) the problem and motivation, (2) your proposed approach, (3) key experimental results, and (4) main conclusions. Keep it to one paragraph.]
\end{abstract}

%%%%%%%%%%%%%%%%%%%%%%%%%%%%%%%%%%%%%%%%%%%%%%%%%%%%%%%%%%%%%%%%
%% Main Paper Content
%%%%%%%%%%%%%%%%%%%%%%%%%%%%%%%%%%%%%%%%%%%%%%%%%%%%%%%%%%%%%%%%

%% Introduction
\section{Introduction}
\label{sec:intro}

% Motivation: Why is this problem important?
Federated reinforcement learning (FRL) enables multiple agents to collaboratively learn policies while preserving data locality and privacy. In this paradigm, agents interact with their local environments and periodically share knowledge through a central aggregation mechanism, avoiding the need to transmit raw interaction data. This approach is particularly valuable in domains where data privacy is critical or communication costs are prohibitive, such as mobile robotics, personalized recommendation systems, and healthcare applications.

However, extending federated learning principles to reinforcement learning presents unique challenges:



% Gap: What is missing in current approaches?
Existing federated learning approaches, particularly those employing FedAvg-style aggregation, face significant limitations when applied to deep reinforcement learning. First, deep neural network representations lose expressiveness during the linear averaging process, as the aggregation operates on weight space rather than the learned feature representations. This representation-agnostic aggregation can degrade the quality of the learned Q-functions. Second, standard federated averaging requires all participating agents to maintain identical model architectures, preventing heterogeneous deployment scenarios where agents may have different computational capabilities or environmental complexities.


% Our approach: What do we propose?
To address these limitations, we propose FedQHD, a federated Q-learning framework that employs hyperdimensional computing for function approximation. By representing Q-functions through high-dimensional vector spaces, our approach enables linear aggregation of learned models while preserving the expressiveness of the approximated value functions. The use of hyperdimensional regression provides a principled method for combining decentralized learning experiences through simple averaging operations.




% Contributions: What are the main contributions?
The main contributions of this work are:
\begin{itemize}
    \item \textbf{Aggregation-aware Q-learning}: We demonstrate that hyperdimensional function approximation maintains expressiveness under linear averaging, enabling effective aggregation of decentralized interaction data without degrading the learned representations.
    \item \textbf{Model heterogeneity support}: Our framework naturally accommodates agents with different model configurations, allowing federated learning across heterogeneous computational environments.
\end{itemize}


% Organization: How is the paper structured?
The rest of this paper is organized as follows. Section~\ref{sec:related_work} reviews related work. Section~\ref{sec:preliminaries} provides background on federated learning, Q-learning, and hyperdimensional computing. Section~\ref{sec:methodology} describes our proposed FedQHD approach. Section~\ref{sec:experiments} details our experimental setup. Section~\ref{sec:results} presents our experimental results. Section~\ref{sec:discussion} discusses the implications of our findings. Finally, Section~\ref{sec:conclusion} concludes the paper.


%% Related Work
\section{Related Work and Background}
\label{sec:related_work}


\subsection{Hyperdimensional Computing (HDC)}
\label{subsec:hdc_background}

Hyperdimensional computing (HDC) is a brain-inspired computational paradigm in which symbols and structured entities are represented as high-dimensional vectors---called \emph{hypervectors}---with components drawn independently from simple distributions~\cite{kanerva2009hyperdimensional}. In such spaces, independently sampled hypervectors are nearly orthogonal with high probability, a geometric property that underlies classical HDC operations such as \emph{bundling} (superposition via addition), \emph{binding} (association via elementwise multiplication or permutation), and \emph{permutation} (role shifting). Because information is distributed holographically across all dimensions, HDC representations exhibit strong robustness to noise, quantization, and partial corruption.

Given an input $x$, an HDC encoder produces a bounded hypervector $h=\phi(x)\in\mathbb{R}^D$ (or $\mathbb{C}^D$) using random projections or compositional schemes built from a small set of base hypervectors. These encoders admit a precise random-feature interpretation: the empirical kernel
\begin{equation}
\label{eq:empirical_kernel}
k_D(x,x')=\langle \phi(x),\phi(x')\rangle
\end{equation}
converges to a smooth limiting kernel $k_\ast(x,x')$ as $D\to\infty$, with uniform concentration at rate $O(D^{-1/2})$.

The following concentration property, standard in the random-feature literature, quantifies this convergence.

With probability at least $1-\delta$ over the random generation of hypervectors, the empirical kernel concentrates uniformly around the limiting kernel:
\begin{equation}
\label{eq:kernel_concentration}
\big| k_D(x,x') - k^\ast(x,x') \big| \leq \frac{C_k}{\sqrt{D}},
\quad \forall\, x,x' \in \mathcal{S} \times \mathcal{A},
\end{equation}
where $C_k > 0$ is an absolute constant depending only on the encoding scheme.

This $O(D^{-1/2})$ concentration rate is established for various random-feature constructions, including Random Fourier Features and related schemes~\cite{rahimi2007random,bach2015equivalence,rudi2017generalization}. As a consequence, linear prediction in the hypervector domain serves as a computationally efficient, finite-dimensional approximation to kernel methods in the RKHS associated with $k_\ast$.




\cite{ni2022qhd} proposes a hyperdimensional regression model to approximate the Q-value function.

\cite{ye2023fedfm} anchor based feature matching FRL

\cite{jin2022federated}  Federated Reinforcement Learning with Environment Heterogeneity (FRL-EH) framework

\cite{hwang2025fed} Federated Reinforcement Learning with Environment Heterogeneity (FRL-EH) framework by incorporating a regularization term in the local update

%% Preliminaries / Background
\section{Preliminaries}
\label{sec:preliminaries}

\subsection{Markov decision processes and off-policy value learning}

\begin{equation}
Q^\star(s,a)= r(s,a) + \gamma\,\mathbb{E}_{s'\sim P(\cdot|s,a)}
\left[\max_{a'\in\mathcal{A}} Q^\star(s',a')\right].
\label{eq:bellman}
\end{equation}
Keep the RL background minimal here and focus to have space in the next section to talk about federation

\subsection{Hyperdimensional computing (HDC)}

\label{subsec:hdc_background}
{\color{red} DESCRIBE ONLY GENERAL HDC HERE COMPACT. }
Hyperdimensional computing (HDC) is a brain-inspired computational paradigm in which symbols and structured entities are represented as high-dimensional vectors---called \emph{hypervectors}---with components drawn independently from simple distributions~\cite{kanerva2009hyperdimensional}. In such spaces, independently sampled hypervectors are nearly orthogonal with high probability, a geometric property that underlies classical HDC operations such as \emph{bundling} (superposition via addition), \emph{binding} (association via elementwise multiplication or permutation), and \emph{permutation} (role shifting). Because information is distributed holographically across all dimensions, HDC representations exhibit strong robustness to noise, quantization, and partial corruption.

Given an input $x$, an HDC encoder produces a bounded hypervector $h=\phi(x)\in\mathbb{R}^D$ (or $\mathbb{C}^D$) using random projections or compositional schemes built from a small set of base hypervectors. These encoders admit a precise random-feature interpretation: the empirical kernel
\begin{equation}
\label{eq:empirical_kernel}
k_D(x,x')=\langle \phi(x),\phi(x')\rangle
\end{equation}
converges to a smooth limiting kernel $k_\ast(x,x')$ as $D\to\infty$, with uniform concentration at rate $O(D^{-1/2})$ \cite{rahimi2007random,bach2015equivalence,rudi2017generalization}. As a consequence, linear prediction in the hypervector domain serves as a computationally efficient, finite-dimensional approximation to kernel methods in the RKHS associated with $k_\ast$.



% \begin{enumerate}
%     \item \textbf{Reward heterogeneity}
%     \item \textbf{}
% \end{enumerate}


% \subsection{Environmental and Model-wise Heterogeneity}
% \label{subsec:hetero}
% We consider $N$ agents. Agent $i$ interacts with its own environment modeled by an MDP
% \begin{equation}
% \mathcal{M}^{(i)}=(\mathcal{S},\mathcal{A},P^{(i)},r^{(i)},\gamma),
% \end{equation}
% where all agents share the same state--action space $(\mathcal{S},\mathcal{A})$ but may differ in reward and transition dynamics.
% For notational convenience, we write $x=(s,a)\in\mathcal{S}\times\mathcal{A}$.

% We quantify environmental mismatch across agents using
% \begin{equation}
% \label{def:hetero}
% \begin{aligned}
% \varepsilon_r
% &:=
% \max_{i,j}\ \frac{\sup_{(s,a)}\big|r^{(i)}(s,a)-r^{(j)}(s,a)\big|}{R},
% \,\\
% \varepsilon_p
% &:=
% \max_{i,j}\ \sup_{(s,a)}
% d_{\mathrm{TV}}\!\big(P^{(i)}(\cdot\mid s,a),P^{(j)}(\cdot\mid s,a)\big), \\
% \varepsilon_m
% &:=
% \max_{i,j}\ 

% \end{aligned}
% \end{equation}
% where $d_{\mathrm{TV}}$ denotes total variation distance.\footnote{Other equivalent discrepancy measures can be used; we adopt total variation for concreteness.}, R denotes the 
% In the finite-time analysis, heterogeneity enters through an induced term $\Lambda(\varepsilon_p,\varepsilon_r)$, which captures the irreducible bias floor arising from aggregating update directions generated by non-identical environments.








%% Methodology (Proposed Approach)



\section{Problem Formulation}
\label{subsec:problem_formulation}
\subsection{QHD: hyperdimensional Q-learning for continuous-state control}
We now describe how hyperdimensional representations are used to approximate action-value functions in continuous-state, discrete-action control. Agent $i$ selects an HDC encoder
\begin{equation}
\Phi_i:\mathcal{S}\to\mathbb{K}^{D_i},
\qquad \mathbb{K}\in\{\mathbb{R},\mathbb{C}\},
\label{eq:phi_i}
\end{equation}
which maps continuous states $s\in\mathbb{R}^{d}$ to hypervectors. For each action $a\in\mathcal{A}$, the agent maintains a parameter hypervector $\mathbf{w}_{i,a}\in\mathbb{K}^{D_i}$ and defines the Q-function approximation
\begin{equation}
Q_i(s,a) \;\triangleq\; \Re\!\left(\langle \Phi_i(s),\,\mathbf{w}_{i,a}\rangle\right).
\label{eq:qhd_q}
\end{equation}
Stacking parameters across actions yields $W_i [\mathbf{w}_{i,a}]_{a\in\mathcal{A}}\in\mathbb{K}^{D_i\times|\mathcal{A}|}$ and $Q_i(s,\cdot)=\Re(\Phi_i(s)^{\mathrm{H}} W_i)$. Although the encoder $\Phi_i(\cdot)$ may be highly nonlinear in $s$,
the model is \emph{linear in parameters}. This linearity is the key structural property we will exploit for federated learning.

Given a replayed transition $(s,a,r,s')$ and a stabilized target estimate (e.g., Double Q-learning with a delayed target copy $Q_i^{-}$), define
\begin{equation}
y = r + \gamma\, Q_i^{-}\!\left(s',\,\arg\max_{a'\in\mathcal{A}} Q_i(s',a')\right).
\label{eq:target}
\end{equation}
QHD performs a lightweight rank-one update
\begin{equation}
\mathbf{w}_{i,a} \leftarrow \mathbf{w}_{i,a} + \eta\,\big(y - Q_i(s,a)\big)\,\Phi_i(s),
\label{eq:qhd_update}
\end{equation}
where $\eta>0$ is a step size...

{\color{red} explain the updates}





% \subsection{HDC Q-function Representation}
% \label{subsec:hdc_q_representation}

% Hyperdimensional computing (HDC) encodes symbols using high-dimensional randomized
% hypervectors with near-orthogonality properties, enabling rich compositional structure while
% preserving simple arithmetic operations. In our setting, an HDC encoder
% \(
%     \Phi: \mathcal{S}\times\mathcal{A}\to\mathbb{C}^D
% \)
% maps each state--action pair $x=(s,a)$ to a $D$-dimensional hypervector $\Phi(x)$.
% The encoder is sampled once at initialization and shared across all agents.
% To ensure numerical stability and compatibility with our finite-time analysis, we assume
% \(
%     \|\Phi(s,a)\|_2 \le 1
%     \,\text{for all $(s,a)$}.
% \)
% Each agent $i$ maintains a parameter hypervector $\mathbf{M}^{(i)}\in\mathbb{C}^D$ and evaluates its
% action--value function as
% \begin{equation}
%     Q_{\mathbf{M}}(s,a)
%     \;=\;
%     \langle \Phi(s,a),\, \mathbf{M} \rangle,
%     \label{eq:hdc_q_function}
% \end{equation}
% where $\langle\cdot,\cdot\rangle$ is the standard inner product on $\mathbb{C}^D$.
% This parameterization exhibits three key structural properties:

% \textit{1)Linearity in parameters:}
%     The map $M \mapsto Q_M$ is linear, so the HDC semi-gradient update 
%     reduces to a rank-one adjustment with $O(D)$ cost, matching the 
%     computational footprint of linear federated TD methods.

% \textit{2)Nonlinearity in inputs:}
%     The encoder $\Phi(s,a)$ is generally nonlinear in $(s,a)$, enabling 
%     $Q_M$ to capture dependencies that cannot be expressed by linear 
%     features. As a result, the HDC representation behaves like a 
%     high-capacity function class while retaining computational simplicity.

% \textit{3)Random-feature kernel structure:}
%     The induced kernel $k_D(x,x')=\langle \Phi(x),\Phi(x')\rangle$ 
%     concentrates to a smooth limiting kernel $k_\ast$ as $D\to\infty$. 
%     Thus, $Q_M$ can be viewed as a finite-dimensional approximation to an 
%     RKHS value function.




\subsection{Federated value learning as a global objective}
We consider $N$ agents learning concurrently, each interacting with its own environment and collecting a private replay buffer $\mathcal{D}_i$. ... 


Federated reinforcement learning seeks to aggregate local value knowledge into a global representation without centralizing data, data sharing is impossible due to privacy, communication limitations, latency, etc. 




\paragraph{Federation as function-space averaging (synchronization operator).}
Each client $i$ learns a local action-value model $Q_i$ from private off-policy replay using QHD-style updates. The server does not access client data or optimize client Bellman losses; instead, it periodically performs a \emph{synchronization} step that aggregates the learned client models into a global teacher in \emph{function space}. Let $\mu$ denote an agreement/query distribution over state--action pairs (instantiated in our method by a small shared anchor set), and let $\pi_i\ge 0$ be federation weights with $\sum_{i=1}^N \pi_i=1$ (e.g., proportional to client data volume).
Given fixed client models $\{Q_i\}_{i=1}^N$, the server defines the global teacher $Q^{\mathrm{glob}}$ as the unique minimizer of
\begin{equation}
Q^{\mathrm{glob}} \;\in\; \arg\min_{Q} \sum_{i=1}^N \pi_i\, E_{(s,a)\sim\mu}\!\left[\big(Q_i(s,a)-Q(s,a)\big)^2\right].
\label{eq:fed_sync}
\end{equation}
This objective is to combine the clients' knowledge in a shared output space rather than in an (often incompatible) parameter space.  The minimizer is the weighted pointwise average (in $\mu$-expectation),
\begin{equation}
Q^{\mathrm{glob}}(s,a) \;=\; \sum_{i=1}^N \pi_i\,Q_i(s,a) \qquad (\mu\text{-a.e.}).
\label{eq:qglob_avg}
\end{equation}

{\color{red} introduce the general problem of function space Q Learning, which can be implied to tabular or continuous state-action setting, to linear or non-linear paprameter averaging in reinforcement learning methond. claim the obstacless}


For TABLUAR, THE solution on top is like hsaring their Q-table, which leads to the exact function-sapce soltuions knowln asn Federated Q-Learning.  FEDERATED AVERAGING HOLDS THE EXACT FUNCTION SPACE, SIMILAR TO A LINEAR FORM LIKE QHD. BUT, FOR NONLINEAR, THEY RELY ON THE HEURISTIC PARAMETER AVERAGING. 



Existing federated RL systems face two central obstacles:{(1) No closed-form, fast synchronization under nonlinear models.} For nonlinear approximators, federated aggregation generally cannot be computed in closed form; client-side assimilation requires iterative optimization, which is expensive and unstable under moving RL targets. {(2) Model heterogeneity breaks parameter aggregation.} Classical federated methods assume a shared parameterization. If clients use different encoders, different architectures, or different parameter dimensions $D_i$, parameter averaging is undefined or semantically meaningless. Such heterogeneity is unavoidable in real autonomous systems.



\section{FedQHD: Function-Space Federation with Closed-Form Translation}
\label{sec:methodology}
The federation objective in Eq.~\eqref{eq:fed_sync} defines a synchronization operator that aggregates client knowledge in \emph{function space}, independently of how local value functions are parameterized or learned. In this section, we show that QHD provides a uniquely favorable structure for realizing this operator efficiently and exactly.

The key enabling property is that, although QHD induces expressive nonlinear approximation over continuous state spaces through the encoder $\Phi_i(\cdot)$, the resulting Q-functions are \emph{linear in parameters}. This linearity allows the federation objective to be solved in closed form when representations are shared, and to be extended to heterogeneous encoders via closed-form translation operators. We first analyze the homogeneous-encoder regime, where the federation operator admits an exact solution that coincides with parameter averaging. This case clarifies the relationship between FedQHD and FedAvg and motivates the general heterogeneous construction that follows.

\subsection{Fed-QHD with Homogeneous Encoders}
\label{sec:fedqhd_homogeneous}

\begin{equation}
Q^{\mathrm{glob}}(s,a) = \langle \Phi(s), w^{\mathrm{glob}}_a \rangle,
\qquad
W^{\mathrm{glob}} \in \mathbb{K}^{D \times |\mathcal{A}|}.
\label{eq:qhd_global}
\end{equation}
Substituting \eqref{eq:qhd_local} and \eqref{eq:qhd_global} into \eqref{eq:fed_sync_recall} yields the quadratic objective
\begin{equation}
\min_{W^{\mathrm{glob}}}\sum_{i=1}^N \pi_i E_{(s,a)\sim\mu} \big[ \langle \Phi(s), (W_i - W^{\mathrm{glob}})_{:,a} \rangle^2\big].
\label{eq:fed_quad}
\end{equation}
Let $\Sigma_\mu = \E_{(s,a)\sim\mu}[\Phi(s)\Phi(s)^\mathrm{H}]$ denote the
feature covariance induced by the agreement distribution.
Then \eqref{eq:fed_quad} can be written as
\begin{equation}
\min_{W^{\mathrm{glob}}} \sum_{i=1}^N \pi_i \sum_{a\in\mathcal{A}} (W_{i,a}-W^{\mathrm{glob}}_a)^\mathrm{H} \Sigma_\mu (W_{i,a}-W^{\mathrm{glob}}_a),
\label{eq:fed_quad_cov}
\end{equation}
which is a strictly convex quadratic problem whenever $\Sigma_\mu \succ 0$, with solution as:
\begin{equation}
W^{\mathrm{glob}} = \sum_{i=1}^N \pi_i W_i.
\label{eq:exact_avg}
\end{equation}

{\color{red} remind that $\Sigma_\mu$ doesn't explicitly appear in solution but involves the raw data sharing with homo encoders}


\subsection{Fed-QHD with Heterogeneous Encoders}
\label{sec:fedqhd_heterogeneous}
Consider a case where each client may use a different hyperdimensional encoder, different kernel hyperparameters, and even a different feature dimension.  In this setting, parameter averaging is undefined, but function-space consensus remains meaningful. Client $i$ uses an encoder \(\Phi_i:\mathcal{S}\to\mathbb{K}^{D_i}, \mathbb{K}\in\{\mathbb{R},\mathbb{C}\},\)
where both the dimension $D_i$ and encoder design (random seed, binding/permutation rules, kernel scale, quantization)
may differ across clients. 

To define federation operationally, we instantiate the agreement distribution $\mu$ by a finite \emph{server-defined anchor set}. Let $\mathcal{S}_{\mathrm{ref}}=\{s_1,\dots,s_m\}$ denote a fixed set of reference states generated independently of all clients. In many autonomy domains, the server can obtain $\mathcal{S}_{\mathrm{ref}}$ from a simulator/digital twin by executing a random or exploratory policy and logging visited states without rewards, actions, or trajectories.

Each client computes its anchor predictions $Q^{\mathrm{ref}}_i[j,a] \;\defeq\; Q_i(s_j,a),\, Q^{\mathrm{ref}}_i\in\mathbb{R}^{m\times|\mathcal{A}|},$
and the server forms the function-space consensus teacher
$Q^{\mathrm{glob}}_{\mathrm{ref}}=\sum_{i=1}^N \pi_i\, Q^{\mathrm{ref}}_i, \sum_{i=1}^N \pi_i=1,\;\pi_i\ge 0$,
which is the unique minimizer of the anchor consensus objective
$\min_{Q\in\mathbb{R}^{m\times|\mathcal{A}|}} \sum_i \pi_i \|Q^{\mathrm{ref}}_i - Q\|_F^2$.


DEFINE AN ANCHOR FOR CLIENT I

For each client $i$, define the \emph{anchor feature matrix} by stacking encoded anchor states:
\begin{equation}
X_i \;=\; 
\begin{bmatrix}
\Phi_i(s_1) \\
\Phi_i(s_2) \\
\vdots \\
\Phi_i(s_m)
\end{bmatrix}
\in \mathbb{K}^{m \times D_i}.
\label{eq:anchor_matrix_i}
\end{equation}
The local Q-function evaluations on the anchor set can then be written compactly as
\begin{equation}
Q^{\mathrm{ref}}_i = \Re(X_i W_i),
\qquad W_i = [\mathbf{w}_{i,a}]_{a\in\mathcal{A}} \in \mathbb{K}^{D_i \times |\mathcal{A}|}.
\label{eq:anchor_q_matrix_form}
\end{equation}

THEN WRITE THE REGULARIZED NONLINEAR EQUATIONS

We seek a client-side parameter matrix $W^{\mathrm{glob}}_i\in\mathbb{K}^{D_i\times|\mathcal{A}|}$ that enables client $i$ to reproduce the global teacher $Q^{\mathrm{glob}}_{\mathrm{ref}}$ as closely as possible using its own encoder $\Phi_i$. This leads to the \emph{regularized least-squares compilation problem}:
\begin{equation}
W^{\mathrm{glob}}_i \;=\; \arg\min_{W \in \mathbb{K}^{D_i \times |\mathcal{A}|}} 
\left\|X_i W - Q^{\mathrm{glob}}_{\mathrm{ref}}\right\|_F^2 
+ \lambda \|W\|_F^2,
\label{eq:regularized_compilation}
\end{equation}
where $\lambda \ge 0$ is a regularization parameter that controls model complexity and stabilizes the solution when the anchor set is small ($m < D_i$) or when $X_i$ is ill-conditioned.


AND THE CONVEX SOLUTIUON FOR IT. 

Problem~\eqref{eq:regularized_compilation} is a convex quadratic optimization problem in $W$ with a unique global minimizer. Taking the gradient with respect to $W$ and setting it to zero yields the regularized normal equations:
\begin{equation}
\left(X_i^\mathrm{H} X_i + \lambda I_{D_i}\right) W^{\mathrm{glob}}_i = X_i^\mathrm{H} Q^{\mathrm{glob}}_{\mathrm{ref}}.
\label{eq:normal_equations}
\end{equation}
When $\lambda > 0$ or when $X_i^\mathrm{H} X_i$ is full rank, the system~\eqref{eq:normal_equations} is invertible, yielding the explicit solution:
\begin{equation}
W^{\mathrm{glob}}_i = \left(X_i^\mathrm{H} X_i + \lambda I_{D_i}\right)^{-1} X_i^\mathrm{H} Q^{\mathrm{glob}}_{\mathrm{ref}}.
\label{eq:ridge_solution}
\end{equation}
In the unregularized limit $\lambda \to 0$, this reduces to the Moore-Penrose pseudoinverse solution:
\begin{equation}
W^{\mathrm{glob}}_i = X_i^\dagger Q^{\mathrm{glob}}_{\mathrm{ref}},
\qquad X_i^\dagger = (X_i^\mathrm{H} X_i)^{-1} X_i^\mathrm{H},
\label{eq:pseudoinverse_solution}
\end{equation}
provided $\mathrm{rank}(X_i) = \min(m, D_i)$.



SEE IF YOU CAN CONNECT HERE WITH ONE EQUATION TO GRAM MATRIX


PROCESS: CLIENT SHARE THEIR DATA, THE SERVER COMPUTES THE Q GLOBAL, AND RETURN THE .... 



% While the shared encoder approach provides efficient aggregation through direct parameter
% averaging, practical federated deployments often involve agents with heterogeneous
% computational capabilities or environmental complexities. To support such scenarios, we
% extend FedQHD to accommodate \textit{agent-specific encoders} $\{\Phi^{(i)}\}_{i=1}^N$
% while preserving the ability to aggregate learned knowledge.

% \subsubsection{Motivation for Heterogeneous Encoders}

% Different agents may benefit from different encoder architectures for several reasons:
% \begin{itemize}
%     \item \textbf{Computational constraints:} Edge devices may require lower-dimensional
%     encoders than datacenter agents.
%     \item \textbf{Environmental complexity:} Agents operating in complex environments may
%     need richer feature representations.
%     \item \textbf{Local adaptation:} Fine-tuning encoders to local state distributions can
%     improve sample efficiency.
% \end{itemize}

% However, with heterogeneous encoders, direct parameter averaging is no longer valid since
% $\mathbf{M}^{(i)}$ lives in the feature space induced by $\Phi^{(i)}$. We address this
% challenge through \textit{regression-based aggregation}.









\section{Projection Residual as a Structural Heterogeneity Measure}
\label{sec:projection_residual}

We formalize the least-squares compilation residual induced by heterogeneous client representations and show that it admits an exact operator form with a clear geometric and optimization-theoretic interpretation. This residual isolates irreducible representation mismatch across clients and provides a principled diagnostic for capacity-limited transfer in heterogeneous federated reinforcement learning.


\begin{proposition}[Irreducible anchor mismatch]
\label{prop:irreducible}
Assume $\lambda=0$ and let $P_i$ denote the orthogonal projector onto
$\mathrm{col}(X_i)\subseteq\mathbb{K}^m$. Then
\[
\min_{W\in\mathbb{K}^{D_i\times|A|}}\|X_iW-Q^{\mathrm{glob}}_{\mathrm{ref}}\|_F
\;=\;\|(I-P_i)Q^{\mathrm{glob}}_{\mathrm{ref}}\|_F.
\]
In particular, $\|(I-P_i)Q^{\mathrm{glob}}_{\mathrm{ref}}\|_F=0$ if and only if
$Q^{\mathrm{glob}}_{\mathrm{ref}}\in\mathrm{col}(X_i)^{|A|}$, i.e., the teacher is
exactly representable by client $i$ on the anchor set.
\end{proposition}


\begin{proposition}[Residual amplification via subspace misalignment]
\label{prop:heterogeneity}
Assume $\lambda=0$ and let $P_i$ denote the orthogonal projector onto
$\mathrm{col}(X_i)$. Suppose the anchor teacher is formed as a convex
combination
\(
Q^{\mathrm{glob}}_{\mathrm{ref}}
=
\sum_{j=1}^N \pi_j Q^{(j)}_{\mathrm{ref}},
\qquad
Q^{(j)}_{\mathrm{ref}}\in\mathrm{col}(X_j),
\quad
\sum_j\pi_j=1.
\)
Then the irreducible mismatch for client $i$ satisfies
\[
R_{i,0}= (I-P_i)Q^{\mathrm{glob}}_{\mathrm{ref}} = \sum_{j\neq i}\pi_j (I-P_i)Q^{(j)}_{\mathrm{ref}}, \text{ with the bound } \|R_{i,0}\|_F \;\le\; \sum_{j\neq i}\pi_j \sin(\theta_{ij}) \,\|Q^{(j)}_{\mathrm{ref}}\|_F,
\]
where $\theta_{ij}$ is the largest principal angle between the subspaces
$\mathrm{col}(X_i)$ and $\mathrm{col}(X_j)$.
\end{proposition}

% \subsection{Federated Learning Protocol}
% \label{subsec:federated_protocol}

% With a shared encoder $\Phi$ across all agents, the linearity of HDC Q-functions in the
% parameter space enables a simple and efficient federated aggregation strategy: direct
% parameter averaging. This approach leverages the parameter-function equivalence property
% (Theorem~\ref{thm:parameter_function_equivalence}), which guarantees that averaging
% parameter vectors is equivalent to averaging Q-functions pointwise.

% \subsection{Direct Parameter Averaging}
% \label{subsec:federated_averaging}
% FedQHD operates in synchronous communication rounds indexed by $k=0,1,2,\dots$. Each
% round consists of three phases with $K$ local update steps, so the global iteration count and
% round index are related by $t = kK + u$, where $u \in \{0,\dots,K-1\}$ denotes
% the local step within round $k$.

% \subsubsection{Broadcast}
% The server maintains a global hypervector $\mathbf{M}_k$ and broadcasts it along with the
% shared encoder $\Phi$ to all clients. Each client initializes its local parameter to $\mathbf{M}_k$.

% \subsubsection{Local Updates}
% Agent $i$ executes $K$ on-policy HDC TD updates using the shared encoder $\Phi$,
% generating
% \begin{equation}
% \begin{split}
%     \mathbf{M}^{(i)}_{k,0}=\mathbf{M}_k,
%     \quad
%      \mathbf{M}^{(i)}_{k,u+1}
%     =
%     \mathbf{M}^{(i)}_{k,u}
%     + \alpha_{k,u}\, \delta_t^{(i)} \Phi(s_t^{(i)}, a_t^{(i)})
% \end{split}
%     \label{eq:local_update}
% \end{equation}
% for $u=0,\dots,K-1$, where $\delta_t^{(i)}$ is the TD error computed using the local Q-function
% $Q_{\mathbf{M}^{(i)}_{k,u}}(s,a) = \langle \Phi(s,a), \mathbf{M}^{(i)}_{k,u} \rangle$.

% \subsubsection{Upload and Aggregation}
% The server receives the local parameter vectors $\{\mathbf{M}^{(i)}_{k,K}\}_{i=1}^N$ and performs
% direct federated averaging:
% \begin{equation}
%     \mathbf{M}_{k+1}
%     =
%     \frac{1}{N}\sum_{i=1}^{N}
%     \mathbf{M}^{(i)}_{k,K}.
%     \label{eq:federated_averaging}
% \end{equation}
% By Theorem~\ref{thm:parameter_function_equivalence}, this parameter averaging is equivalent
% to averaging the Q-functions, i.e., $Q_{\mathbf{M}_{k+1}}(s,a) = \frac{1}{N}\sum_{i=1}^N Q_{\mathbf{M}^{(i)}_{k,K}}(s,a)$
% for all $(s,a)$. The averaged model $\mathbf{M}_{k+1}$ is then broadcast to all clients to
% initiate round $k+1$.
\subsection{FedQHD Algorithm}
\label{subsec:algorithm}

% We present the complete FedQHD algorithm with shared encoder and direct parameter averaging.
% This approach leverages the parameter-function equivalence property to enable efficient
% federated aggregation with minimal communication overhead ($O(D)$ parameter uploads per round).

\begin{algorithm}[t]
\caption{FedQHD: Shared Encoder Protocol}
\label{alg:fedqhd_shared}
\begin{algorithmic}[1]
\STATE \textbf{Input:} Rounds $T$, local steps $K$, agents $N$, learning rates $\{\alpha_{k,u}\}$
\STATE \textbf{Server Initialize:}
\STATE \quad Shared encoder $\Phi: \mathcal{S}\times\mathcal{A}\to\mathbb{C}^D$
\STATE \quad Global model $\mathbf{M}_0 = \mathbf{0} \in \mathbb{C}^D$
\STATE
\FOR{round $k = 0, 1, \ldots, T-1$}
    \STATE \textbf{[Server Broadcast]}
    \STATE \quad Send $\mathbf{M}_k$ to all agents (encoder $\Phi$ sent once at initialization)
    \STATE
    \STATE \textbf{[Client Local Learning]} (in parallel for all agents $i = 1, \ldots, N$)
    \STATE \quad Initialize: $\mathbf{M}^{(i)}_{k,0} = \mathbf{M}_k$
    \FOR{local step $u = 0, 1, \ldots, K-1$}
        \STATE \quad Observe state $s_t^{(i)}$, select action $a_t^{(i)}$ via $\epsilon$-greedy on $Q_{\mathbf{M}^{(i)}_{k,u}}$
        \STATE \quad Execute $a_t^{(i)}$, observe reward $r_t^{(i)}$ and next state $s_{t+1}^{(i)}$
        \STATE \quad Select next action $a_{t+1}^{(i)}$ using current policy (SARSA)
        \STATE \quad Encode: $\phi_t = \Phi(s_t^{(i)}, a_t^{(i)})$, $\phi_{t+1} = \Phi(s_{t+1}^{(i)}, a_{t+1}^{(i)})$
        \STATE \quad TD error: $\delta_t^{(i)} = r_t^{(i)} + \gamma \langle \phi_{t+1}, \mathbf{M}^{(i)}_{k,u} \rangle - \langle \phi_t, \mathbf{M}^{(i)}_{k,u} \rangle$
        \STATE \quad \textbf{Update:} $\mathbf{M}^{(i)}_{k,u+1} = \mathbf{M}^{(i)}_{k,u} + \alpha_{k,u} \delta_t^{(i)} \phi_t$
    \ENDFOR
    \STATE \quad \textbf{Upload:} Send $\mathbf{M}^{(i)}_{k,K}$ to server
    \STATE
    \STATE \textbf{[Server Aggregation]}
    \STATE \quad Receive parameter vectors $\{\mathbf{M}^{(i)}_{k,K}\}_{i=1}^{N}$ from all agents
    \STATE \quad \textbf{Average parameters:} $\mathbf{M}_{k+1} = \frac{1}{N}\sum_{i=1}^{N} \mathbf{M}^{(i)}_{k,K}$
    \STATE \quad \textbf{(Optional)} Project: $\mathbf{M}_{k+1} \leftarrow \text{Proj}_{\mathcal{B}(G)}(\mathbf{M}_{k+1})$
\ENDFOR
\STATE
\STATE \textbf{Output:} Global Q-function $Q_{\mathbf{M}_T}(s,a) = \langle \Phi(s,a), \mathbf{M}_T \rangle$
\end{algorithmic}
\end{algorithm}

% \subsection{Theoretical Property: Parameter-Function Equivalence}
% \label{subsec:theory}

% \begin{theorem}[Parameter-Function Equivalence]
% \label{thm:parameter_function_equivalence}
% Let $\Phi: \mathcal{S}\times\mathcal{A} \to \mathbb{C}^D$ be a shared encoder fixed across all agents, and let $\{\mathbf{M}^{(i)}\}_{i=1}^N$ be local model hypervectors. Define the averaged hypervector as
% \begin{equation}
%     \bar{\mathbf{M}} = \frac{1}{N}\sum_{i=1}^N \mathbf{M}^{(i)}.
% \end{equation}
% Then, for all state-action pairs $(s,a) \in \mathcal{S} \times \mathcal{A}$, parameter averaging is equivalent to function averaging:
% \begin{equation}
%     Q_{\bar{\mathbf{M}}}(s,a) = \frac{1}{N}\sum_{i=1}^N Q_{\mathbf{M}^{(i)}}(s,a).
%     \label{eq:param_func_equiv}
% \end{equation}
% \end{theorem}

% \begin{proof}
% By definition of the HDC Q-function~\eqref{eq:hdc_q_function} and linearity of the inner product:
% \begin{align}
%     Q_{\bar{\mathbf{M}}}(s,a)
%     &= \langle \Phi(s,a), \bar{\mathbf{M}} \rangle \\
%     &= \left\langle \Phi(s,a), \frac{1}{N}\sum_{i=1}^N \mathbf{M}^{(i)} \right\rangle \\
%     &= \frac{1}{N}\sum_{i=1}^N \langle \Phi(s,a), \mathbf{M}^{(i)} \rangle \quad \text{(linearity of $\langle\cdot,\cdot\rangle$)} \\
%     &= \frac{1}{N}\sum_{i=1}^N Q_{\mathbf{M}^{(i)}}(s,a). \qquad \qed
% \end{align}
% \end{proof}

% This theorem justifies the direct parameter averaging strategy in Algorithm~\ref{alg:fedqhd_shared}
% and distinguishes HDC-based Q-learning from deep neural network approaches, where parameter
% averaging does not preserve function-level aggregation due to nonlinear activation functions.

% \subsection{Extension: Heterogeneous Encoders via Regression-Based Aggregation}
% \label{subsec:heterogeneous_extension}


\subsubsection{Regression-Based Aggregation Protocol}

The key insight is that while parameters cannot be directly averaged across different
feature spaces, Q-\textit{functions} can be averaged pointwise. By exploiting the linearity
of HDC Q-functions in parameter space, we can solve a regression problem to find global
parameters that approximate the averaged Q-function.

\textbf{Anchor-based aggregation:} The server maintains:
\begin{itemize}
    \item A global encoder $\Phi_{\text{global}}: \mathcal{S}\times\mathcal{A}\to\mathbb{C}^D$
    \item A set of anchor state-action pairs $\mathcal{X} = \{(s_j, a_j)\}_{j=1}^{L}$
    sampled from a representative distribution
\end{itemize}

\textbf{Upload phase:} Instead of uploading parameter vectors, each agent evaluates its
local Q-function on the anchor set:
\begin{equation}
    \mathbf{q}^{(i)}_k = \left[ Q_{\mathbf{M}^{(i)}_{k,K}}^{(i)}(s_j, a_j) \right]_{j=1}^{L}
    = \left[ \langle \Phi^{(i)}(s_j, a_j), \mathbf{M}^{(i)}_{k,K} \rangle \right]_{j=1}^{L}
    \in \mathbb{R}^L.
    \label{eq:q_value_upload_het}
\end{equation}
This requires $O(L)$ communication per agent, which can be more efficient than $O(D)$
parameter uploads when $L \ll D$.

\textbf{Aggregation phase:} The server computes the target averaged Q-values:
\begin{equation}
    \bar{\mathbf{q}}_k = \frac{1}{N}\sum_{i=1}^{N} \mathbf{q}^{(i)}_k,
    \label{eq:averaged_q_het}
\end{equation}
and solves a least-squares regression to find global parameters that best approximate
$\bar{\mathbf{q}}_k$ under $\Phi_{\text{global}}$:
\begin{equation}
    \mathbf{M}_{k+1} = \arg\min_{\mathbf{M} \in \mathbb{C}^D}
    \sum_{j=1}^{L} \left\|
    \langle \Phi_{\text{global}}(s_j, a_j), \mathbf{M} \rangle - \bar{q}_{k,j}
    \right\|^2.
    \label{eq:regression_aggregation_het}
\end{equation}

This has a closed-form solution via the Moore-Penrose pseudoinverse:
\begin{equation}
    \mathbf{M}_{k+1} = \mathbf{\Psi}^\dagger \bar{\mathbf{q}}_k,
    \label{eq:closed_form_aggregation_het}
\end{equation}
where $\mathbf{\Psi} = [\Phi_{\text{global}}(s_1,a_1), \dots, \Phi_{\text{global}}(s_L,a_L)]^T \in \mathbb{C}^{L \times D}$.

\subsubsection{Modified Algorithm for Heterogeneous Encoders}

Algorithm~\ref{alg:fedqhd_heterogeneous} presents the complete protocol with regression-based
aggregation.

\begin{algorithm}[t]
\caption{FedQHD: Heterogeneous Encoder Extension}
\label{alg:fedqhd_heterogeneous}
\begin{algorithmic}[1]
\STATE \textbf{Input:} Rounds $T$, local steps $K$, agents $N$, anchor size $L$, learning rates $\{\alpha_{k,u}\}$
\STATE \textbf{Server Initialize:}
\STATE \quad Global encoder $\Phi_{\text{global}}: \mathcal{S}\times\mathcal{A}\to\mathbb{C}^D$
\STATE \quad Global model $\mathbf{M}_0 = \mathbf{0} \in \mathbb{C}^D$
\STATE \quad Anchor set $\mathcal{X} = \{(s_j, a_j)\}_{j=1}^{L}$
\STATE \quad Feature matrix $\mathbf{\Psi} = [\Phi_{\text{global}}(s_1,a_1), \dots, \Phi_{\text{global}}(s_L,a_L)]^T \in \mathbb{C}^{L \times D}$
\STATE \textbf{Client Initialize:} Each agent $i$ initializes or adapts local encoder $\Phi^{(i)}$
\STATE
\FOR{round $k = 0, 1, \ldots, T-1$}
    \STATE \textbf{[Server Broadcast]} Send $(\mathbf{M}_k, \Phi_{\text{global}}, \mathcal{X})$ to all agents
    \STATE
    \STATE \textbf{[Client Local Learning]} (in parallel for all agents $i = 1, \ldots, N$)
    \STATE \quad Initialize: $\mathbf{M}^{(i)}_{k,0} = \mathbf{M}_k$
    \STATE \quad \textbf{(Optional) Adapt encoder:} Fine-tune $\Phi^{(i)}$ based on $\Phi_{\text{global}}$
    \FOR{local step $u = 0, 1, \ldots, K-1$}
        \STATE \quad Observe state $s_t^{(i)}$, select action $a_t^{(i)}$ via $\epsilon$-greedy
        \STATE \quad Execute $a_t^{(i)}$, observe $r_t^{(i)}, s_{t+1}^{(i)}$, select next action $a_{t+1}^{(i)}$
        \STATE \quad Encode: $\phi_t = \Phi^{(i)}(s_t^{(i)}, a_t^{(i)})$, $\phi_{t+1} = \Phi^{(i)}(s_{t+1}^{(i)}, a_{t+1}^{(i)})$
        \STATE \quad TD error: $\delta_t^{(i)} = r_t^{(i)} + \gamma \langle \phi_{t+1}, \mathbf{M}^{(i)}_{k,u} \rangle - \langle \phi_t, \mathbf{M}^{(i)}_{k,u} \rangle$
        \STATE \quad \textbf{Update:} $\mathbf{M}^{(i)}_{k,u+1} = \mathbf{M}^{(i)}_{k,u} + \alpha_{k,u} \delta_t^{(i)} \phi_t$
    \ENDFOR
    \STATE \quad \textbf{Evaluate on anchors:} $\mathbf{q}^{(i)}_k = [\langle \Phi^{(i)}(s_j, a_j), \mathbf{M}^{(i)}_{k,K} \rangle]_{j=1}^{L}$
    \STATE \quad \textbf{Upload:} Send Q-value vector $\mathbf{q}^{(i)}_k \in \mathbb{R}^L$ to server
    \STATE
    \STATE \textbf{[Server Aggregation]}
    \STATE \quad Receive $\{\mathbf{q}^{(i)}_k\}_{i=1}^{N}$ from all agents
    \STATE \quad \textbf{Average Q-values:} $\bar{\mathbf{q}}_k = \frac{1}{N}\sum_{i=1}^{N} \mathbf{q}^{(i)}_k$
    \STATE \quad \textbf{Regression:} $\mathbf{M}_{k+1} = \mathbf{\Psi}^\dagger \bar{\mathbf{q}}_k$
    \STATE \quad \textbf{(Optional)} Project: $\mathbf{M}_{k+1} \leftarrow \text{Proj}_{\mathcal{B}(G)}(\mathbf{M}_{k+1})$
\ENDFOR
\STATE
\STATE \textbf{Output:} Global Q-function $Q_{\mathbf{M}_T}(s,a) = \langle \Phi_{\text{global}}(s,a), \mathbf{M}_T \rangle$
\end{algorithmic}
\end{algorithm}

\subsubsection{Theoretical Guarantee}

\begin{theorem}[Regression-Based Aggregation Correctness]
\label{thm:regression_aggregation}
Let $\{\Phi^{(i)}\}_{i=1}^N$ be heterogeneous local encoders and $\{\mathbf{M}^{(i)}\}_{i=1}^N$ be
corresponding parameter vectors. Let $\Phi_{\text{global}}$ be the global encoder and
$\mathcal{X} = \{(s_j, a_j)\}_{j=1}^{L}$ be anchor points. Define the target averaged Q-function as
\begin{equation}
    \bar{Q}(s,a) = \frac{1}{N}\sum_{i=1}^N Q_{\mathbf{M}^{(i)}}^{(i)}(s,a),
\end{equation}
where $Q_{\mathbf{M}^{(i)}}^{(i)}(s,a) = \langle \Phi^{(i)}(s,a), \mathbf{M}^{(i)} \rangle$.
Let $\mathbf{M}_{\text{reg}}$ be the solution to~\eqref{eq:regression_aggregation_het}.
Then, $\mathbf{M}_{\text{reg}}$ minimizes the mean squared error on the anchor set:
\begin{equation}
    \mathbf{M}_{\text{reg}} = \arg\min_{\mathbf{M}} \sum_{j=1}^{L}
    \left| \langle \Phi_{\text{global}}(s_j, a_j), \mathbf{M} \rangle - \bar{Q}(s_j, a_j) \right|^2.
\end{equation}
Moreover, if $L \geq D$ and the anchor features span $\mathbb{C}^D$, then $Q_{\mathbf{M}_{\text{reg}}}$
exactly interpolates $\bar{Q}$ on $\mathcal{X}$.
\end{theorem}

\begin{proof}
The regression objective is a standard least-squares problem with linear model
$\langle \Phi_{\text{global}}(s_j, a_j), \mathbf{M} \rangle$. The closed-form solution
follows from normal equations. When $L \geq D$ and $\text{rank}(\mathbf{\Psi}) = D$,
the pseudoinverse satisfies $\mathbf{\Psi}^\dagger \mathbf{\Psi} = \mathbf{I}_D$, yielding
exact interpolation on $\mathcal{X}$. \qed
\end{proof}

\textbf{Remark:} Theorem~\ref{thm:regression_aggregation} shows that regression-based
aggregation preserves the averaged Q-function on anchor points, enabling heterogeneous
 model aggregation. The choice of anchor distribution and size $L$ controls the
approximation quality and generalization to unseen states. In practice, anchors can be
sampled from a replay buffer or from the empirical state-action distribution observed
during training.



%% Experiments
\section{Experiments}
\label{sec:experiments}

We validate FedQHD across four continuous-control benchmarks under both homogeneous and heterogeneous encoder settings, and address four research questions.
\textbf{(Q1)} \emph{Homogeneous federation.}
Does federated QHD with a shared encoder outperform independent learning and existing federated RL baselines, and how close does it come to the centralised oracle?
\textbf{(Q2)} \emph{Heterogeneous federation.}
Can anchor-based aggregation bridge per-client encoder mismatches, and what fraction of the homogeneous gain is preserved when encoders differ in dimension and bandwidth?
\textbf{(Q3)} \emph{Scalability.}
Does adding more clients yield consistent, monotone gains, and at what federation size do the returns diminish?
\textbf{(Q4)} \emph{Ablation.}
How do encoder dimension $D$ and anchor set size $m$ govern the quality of the anchor-based projection step?

\subsection{Experimental Setup}
\label{subsec:exp_setup}

\paragraph{Environments and Tasks.}
We evaluate FedQHD on four continuous-state control benchmarks from OpenAI Gym~\cite{brockman2016openai}:
\begin{itemize}
    \item \textbf{CartPole-v1}: Balance a pole on a cart by applying left/right force (4D state: position, velocity, angle, angular velocity; 2 discrete actions). Episode ends at 500 steps.
    \item \textbf{Acrobot-v1}: Swing up a two-link underactuated pendulum to a target height (6D state: $\cos/\sin$ of two joint angles and their angular velocities; 3 discrete torques). Higher reward (closer to 0) is better.
    \item \textbf{LunarLander-v2}: Land a spacecraft between two flags using main and side thrusters (8D state: position, velocity, angle, angular velocity, leg contact; 4 discrete actions). Episode solved at reward $\geq 200$.
    \item \textbf{MountainCar-v0}: Drive an underpowered car to a hilltop by building momentum (2D state: position, velocity; 3 discrete actions). Sparse reward: $-1$ per step until goal reached.
\end{itemize}
All $N=5$ clients share the same environment type but run independent instances with different random seeds, creating statistical variability without requiring environmental heterogeneity.

\paragraph{Hyperdimensional Encoders.}
All QHD-based methods use a random Fourier feature (RFF) encoder~\cite{rahimi2007random}:
\begin{equation}
\Phi(s) = \frac{1}{\sqrt{D}}\bigl[\cos(\omega_1^\top s + b_1),\dots,\cos(\omega_D^\top s + b_D)\bigr]^\top,
\end{equation}
where $\omega_j \sim \mathcal{N}(0, \sigma^2 I)$, $b_j \sim \mathrm{Uniform}[0, 2\pi]$, and $\sigma$ is the RFF bandwidth.
For \textbf{Q1 (homogeneous)}, all clients share a single encoder ($D=10{,}000$, $\sigma=\sigma_0$, seed $=42$), so $\Phi_i \equiv \Phi$ and parameter averaging is exact.
For \textbf{Q2 (heterogeneous)}, each client $i$ draws an independent encoder with bandwidth $\sigma_i \sim \mathrm{Uniform}[0.5\sigma_0, 1.5\sigma_0]$ and dimension $D_i \in \{1{,}000, 5{,}000, 10{,}000, 50{,}000\}$ (assigned cyclically across clients). Anchor-based aggregation is required because $\Phi_i \neq \Phi_j$.

\paragraph{Anchor Set Construction.}
For Q2 aggregation, the server constructs an anchor set $\mathcal{S}_\mathrm{ref}$ of $m=200$ states collected via random rollouts in the environment. This simulates a lightweight server-side simulator or held-out validation set with no access to client training data.

\paragraph{Implementation Details.}
All QHD methods use $\epsilon$-greedy exploration with $\epsilon$ initialised at 0.8--1.0 and annealed exponentially to 0.001 per episode. The QHD linear weights are updated via online TD($0$) with learning rate $\eta_\mathrm{QHD}$; DQN baselines use Adam with $\eta_\mathrm{DQN}=10^{-3}$ and a replay buffer of 10,000 transitions. Hyperparameters are tuned per environment via grid search (Section~\ref{subsec:baselines}); the resulting values are summarised in Table~\ref{tab:hparams}. Federation weights are uniform: $\pi_i = 1/N$. Each experiment uses a single run with random seed 42; results are averaged over the last 100 episodes.

\begin{table}[h]
  \centering
  \caption{Per-environment QHD hyperparameters (Q1/Q2 homogeneous). $K$: aggregation interval (episodes); $\eta$: learning rate; $\gamma$: discount; $\epsilon_\delta$: exploration decay per episode.}
  \label{tab:hparams}
  \setlength{\tabcolsep}{5pt}
  \begin{tabular}{lcccc}
  \toprule
  Environment & $\eta$ & $\gamma$ & $\epsilon_\delta$ & $K$ \\
  \midrule
  CartPole-v1    & 0.20 & 0.90 & 0.900 & 50 \\
  Acrobot-v1     & 0.10 & 0.99 & 0.995 & 10 \\
  LunarLander-v2 & 0.20 & 0.99 & 0.990 & 10 \\
  MountainCar-v0 & 0.10 & 0.99 & 0.995 & 50 \\
  \bottomrule
  \end{tabular}
\end{table}

\subsection{Baselines}
\label{subsec:baselines}

We compare FedQHD against six methods spanning the spectrum from no federation to perfect data sharing:

\begin{enumerate}
  \item \textbf{Independent QHD} (lower bound, Q1\,\&\,Q2).
    Each client trains a local QHD agent on its own environment with no communication. Reports the average reward of the $N$ local agents. Establishes the floor that any federated method must exceed.

  \item \textbf{Oracle QHD}$^\dagger$ (upper bound, Q1\,\&\,Q2).
    $N$ QHD agents are each trained on \emph{all} $N$ client environments simultaneously (perfect data pooling), aggregating parameters every episode. Under Q1 all agents share the same encoder (direct averaging); under Q2 each agent uses its client-specific encoder and anchor-based projection is applied at every episode. Provides the best achievable performance under perfect data sharing.

  \item \textbf{Oracle DQN}$^\dagger$ (upper bound, Q1\,\&\,Q2).
    Same as Oracle QHD but uses a DQN agent (2-layer MLP, 128 hidden units, Adam $\eta=10^{-3}$) trained on pooled data from all environments. Serves as a cross-architecture oracle baseline.

  \item \textbf{FedAvg-DQN}~\cite{mcmahan2017communication} (Q1 only).
    Federated deep Q-learning: each client trains a local DQN for $K=10$ episodes, then the server averages all network weights. Applicable only under homogeneous architectures; excluded from Q2 because weight averaging is undefined across different encoder dimensions.

  \item \textbf{Truncate FedAvg-QHD} (Q1\,\&\,Q2).
    A naive heterogeneous baseline: client weight matrices $W_i \in \mathbb{R}^{D_i \times |\mathcal{A}|}$ are truncated to the minimum dimension $D_{\min}$ before averaging, then zero-padded back to each client's original dimension. Under Q1 this reduces to standard FedAvg-QHD; under Q2 it discards high-frequency features, degrading performance.

  \item \textbf{Distillation FedDQN} (Q1\,\&\,Q2).
    Policy distillation across heterogeneous agents: each client trains a local DQN student to match the soft action distribution (temperature $\tau=1$) of a shared teacher DQN that is updated via FedAvg every $K=25$ episodes. The student can have any architecture, making the method applicable under Q2. Requires additional forward passes through the teacher at every aggregation round, increasing wall-clock cost relative to Q1.
\end{enumerate}


\subsection{Evaluation Metrics}
\label{subsec:metrics}

We evaluate methods using:
\begin{itemize}
    \item \textbf{Average Reward}: Mean episodic return across all client environments, measured every 10 communication rounds.
    \item \textbf{Convergence Speed}: Number of environment steps (or communication rounds) to reach 90\% of oracle performance.
    \item \textbf{Projection Residual}: $\|R_{i,0}\|_F = \|(I - P_i)Q^{\mathrm{glob}}_{\mathrm{ref}}\|_F$ (heterogeneous case only).
    \item \textbf{Oracle Q-Approximation Error} (ablation studies): RMSE between the best-in-class ridge-regression estimate and the reference Q-function $Q^*$ on held-out anchor states,
    \begin{equation}
      \label{eq:metric_qaerr}
      \mathrm{err} = \bigl\|\hat{Q}_{\mathrm{eval}} - Q^*_{\mathrm{eval}}\bigr\|_F,
      \quad
      \hat{W} = (X^\top X + \lambda I)^{-1}X^\top Q^*,\quad
      \hat{Q}_{\mathrm{eval}} = X_{\mathrm{eval}}\hat{W},
    \end{equation}
    where $Q^*$ is approximated via Monte-Carlo rollouts from a converged reference agent ($D_{\mathrm{ref}}=2048$, $V(\pi_{\mathrm{ref}})=500$, 30 rollouts per anchor state on CartPole-v1). This metric isolates approximation capacity from RL training dynamics.
\end{itemize}

\subsection{Results}
\label{sec:results}

\subsubsection{Performance Comparison (Q1 and Q2)}
\label{subsec:performance_homo}

Table~\ref{tab:main_results} reports the final average reward (mean over the last 100 episodes) for all methods across four environments under both encoder settings.
Figure~\ref{fig:learning_curves} shows the smoothed learning curves for CartPole-v1 and LunarLander-v3; results for the remaining environments are consistent and omitted for space.


  \begin{table}[t]
  \centering
  \caption{%
    Final average reward (last 100 episodes, 600-episode runs, $N{=}5$ clients).
    \textbf{Bold}: best non-oracle result per column.
    $\dagger$: oracle upper bounds (centralised data access).
    {--}: method not applicable to this setting.
    Higher is better for all environments.
  }
  \label{tab:main_results}
  \setlength{\tabcolsep}{4pt}
  \begin{tabular}{l rr rr rr rr}
  \toprule
   & \multicolumn{2}{c}{CartPole}
   & \multicolumn{2}{c}{Acrobot}
   & \multicolumn{2}{c}{LunarLander}
   & \multicolumn{2}{c}{MountainCar} \\
  \cmidrule(lr){2-3}\cmidrule(lr){4-5}\cmidrule(lr){6-7}\cmidrule(lr){8-9}
  Method
    & Homo & Hetero
    & Homo & Hetero
    & Homo & Hetero
    & Homo & Hetero \\
  \midrule
  Independent QHD
    &  23.42 &  23.31
    & $-$499.92 & $-$493.93
    & $-$179.99 & $-$179.21
    & $-$200.00 & $-$200.00 \\
  Truncate FedAvg-QHD
    & 193.17 &  51.30
    & $-$103.88 & $-$106.97
    & 70.44 & 0.88
    & $-$142.56 & $-$166.78 \\
  FedAvg-DQN
    & 186.78 & {--}
    & $\mathbf{-84.18}$ & {--}
    & 225.44 & {--}
    & $-$163.55 & {--} \\
  Distillation FedDQN
    & 144.44 & 176.82
    & $-$86.80 & $\mathbf{-88.06}$
    & 116.11 & $\mathbf{131.36}$
    & $-$150.78 & $-$170.54 \\
  FedQHD
    & $\mathbf{456.46}$ & $\mathbf{422.53}$
    & $-$103.79 & $-$101.18
    & $\mathbf{236.26}$ & 101.50
    & $\mathbf{-140.81}$ & $\mathbf{-159.08}$ \\
  \midrule
  Oracle QHD$^\dagger$
    & 498.09 & 500.00
    & $-$99.86 & $-$97.78
    & 83.84 & 232.93
    & $-$138.54 & $-$137.82 \\
  Oracle DQN$^\dagger$
    & 195.12 & 215.48
    & $-$77.16 & $-$80.50
    & 183.45 & 202.56
    & $-$122.45 & $-$119.97 \\
  \bottomrule
  \end{tabular}
  \end{table}


\begin{figure}[t]
  \centering
  \includegraphics[width=\linewidth]{../results/figures/CartPole_LunarLander_combined.pdf}
  \caption{Learning curves (smoothed, window 20) for CartPole-v1 (top) and LunarLander-v3 (bottom) under homogeneous (left) and heterogeneous (right) encoder settings. $N=5$ clients, 600 episodes.}
  \label{fig:learning_curves}
\end{figure}

\paragraph{Q1 — Homogeneous encoders.}
FedQHD achieves the best non-oracle performance in three of four environments.
On CartPole-v1 it reaches 456.5, closing $90\%$ of the gap to Oracle QHD (498.1) while outperforming the next-best federated baseline, FedAvg-DQN (186.8), by a factor of $2.4\times$.
The learning curve (Figure~\ref{fig:learning_curves}a) confirms that FedQHD converges rapidly after episode 200, whereas DQN-based baselines plateau well below 200.
On LunarLander-v3, FedQHD attains 236.3---the only non-oracle method to exceed the solved threshold of 200---while FedAvg-DQN trails at 225.4.
Notably, Oracle QHD scores only 83.8 on LunarLander, \emph{below} FedQHD; this reflects a gradient-conflict effect: aggregating every episode (Oracle QHD's design) destroys the local specialisation that each client develops between communication rounds, a failure mode avoided by FedQHD's $K$-episode aggregation gap.
On MountainCar-v0, FedQHD ($-140.8$) matches Oracle QHD ($-138.5$) almost exactly.
On Acrobot-v1, FedAvg-DQN ($-84.2$) edges FedQHD ($-103.8$), showing that DQN's non-linear function approximation can be advantageous in environments where the linear QHD approximation is under-expressive.

\paragraph{Q2 — Heterogeneous encoders.}
Anchor-based aggregation preserves a substantial fraction of the homogeneous gain on low-dimensional tasks.
CartPole-v1 drops only modestly from 456.5 (Q1) to 422.5 (Q2), retaining $92\%$ of the Q1 performance; the learning curve (Figure~\ref{fig:learning_curves}b) tracks the homogeneous curve closely from episode 300 onward.
The performance gap widens on LunarLander-v3: FedQHD falls to 101.5 (Q2) versus 236.3 (Q1), a $57\%$ reduction.
The Q2 learning curve (Figure~\ref{fig:learning_curves}d) shows that the Oracle QHD hetero ceiling is 232.9, implying that the limitation is not the anchor mechanism per se but the intrinsic difficulty of aligning 8-dimensional encoder outputs across clients with varying bandwidths.
Distillation FedDQN outperforms FedQHD on LunarLander-v3 Q2 (131.4 vs.\ 101.5) and Acrobot-v1 Q2 ($-88.1$ vs.\ $-101.2$), suggesting that knowledge distillation is a stronger signal when encoder heterogeneity is high and the state space is complex.
Truncate FedAvg-QHD degrades sharply under Q2 (CartPole: 193.2 $\to$ 51.3; LunarLander: 70.4 $\to$ 0.9), confirming that naive dimension-matching by truncation destroys the feature structure critical for linear Q-function estimation.

\textbf{Answer to Q1\,\&\,Q2.}
FedQHD is the best non-oracle method in 5 of 8 cells.
Under homogeneous encoders it consistently outperforms all federated DQN variants on CartPole and LunarLander, approaching the centralised oracle.
Anchor-based aggregation successfully extends federation to heterogeneous encoders, though with a performance cost that grows with state-space dimensionality.

\subsubsection{Computation Cost}

Table~\ref{tab:computation_cost} reports wall-clock training time (minutes, 600 episodes, $N=5$ clients; CP = CartPole, Acr = Acrobot, LL = LunarLander, MC = MountainCar).
FedQHD is among the fastest non-trivial methods: under Q1 it takes 9.7--20.3 min across environments, compared to 25.9--85.6 min for FedAvg-DQN and 14.3--69.2 min for Distillation FedDQN.
Under Q2 the advantage is even more pronounced: FedQHD requires only 1.4--5.9 min because anchor-based aggregation replaces neural-network backward passes with a single ridge-regression solve ($\mathcal{O}(D^2|A|)$ per client), cutting training time by $3$--$12\times$ relative to its Q1 cost.
By contrast, Distillation FedDQN's Q2 cost (22.2--71.3 min) matches or exceeds its Q1 cost because policy distillation requires additional forward passes through the teacher network at every aggregation round.
Oracle DQN is the most expensive method overall (up to 156.5 min on LunarLander Q2) due to full DQN training on the pooled dataset across all $N$ environments.

\begin{table}[t]
  \centering
  \caption{%
    Wall-clock training time (minutes, 600 episodes, $N{=}5$ clients).
    Q1 = homogeneous encoders; Q2 = heterogeneous encoders.
    {--}: method not applicable to this setting.
  }
  \label{tab:computation_cost}
  \setlength{\tabcolsep}{4pt}
  \begin{tabular}{l rrrr rrrr}
  \toprule
   & \multicolumn{4}{c}{Q1 (Homogeneous)}
   & \multicolumn{4}{c}{Q2 (Heterogeneous)} \\
  \cmidrule(lr){2-5}\cmidrule(lr){6-9}
  Method & CP & Acr & LL & MC & CP & Acr & LL & MC \\
  \midrule
  Independent QHD
    & 1.7 & 49.4 & 5.6 & 12.7 & 0.1 & 4.5 & 1.1 & 1.4 \\
  Truncate FedAvg-QHD
    & 3.5 & 10.2 & 26.4 & 7.6 & 0.3 & 2.5 & 5.2 & 1.7 \\
  FedAvg-DQN
    & 25.9 & 81.2 & 72.8 & 85.6 & {--} & {--} & {--} & {--} \\
  Distillation FedDQN
    & 14.3 & 46.1 & 69.2 & 28.1 & 22.2 & 47.3 & 71.3 & 52.9 \\
  FedQHD
    & 10.3 & 9.7 & 20.3 & 17.2 & 1.6 & 3.5 & 5.9 & 1.4 \\
  \midrule
  Oracle QHD$^\dagger$
    & 12.3 & 3.9 & 18.9 & 7.7 & 10.2 & 10.9 & 14.8 & 8.5 \\
  Oracle DQN$^\dagger$
    & 24.4 & 67.6 & 63.0 & 67.8 & 81.6 & 86.4 & 156.5 & 101.5 \\
  \bottomrule
  \end{tabular}
  \end{table}



\subsubsection{Scalability Analysis}
\label{subsec:scalability}

We evaluate scalability by varying $N \in \{1, 2, 5, 20\}$ on CartPole-v1 and LunarLander-v3, measuring average reward over the final 30 episodes (Figure~\ref{fig:scalability}).
Following standard benchmarks~\cite{brockman2016openai}, we define a task as \emph{solved} when the rolling average reward exceeds a environment-specific threshold: $\geq\!475$ for CartPole-v1 and $\geq\!200$ for LunarLander-v3.
The dashed red line in each panel marks this threshold; bars that cross it indicate a qualitatively reliable policy.

FedQHD's performance improves consistently with $N$ on both environments, yielding cumulative gains of $+30.6\%$ on CartPole and $+29.5\%$ on LunarLander relative to the single-agent baseline.
The marginal benefit is front-loaded and environment-dependent.
On CartPole, the dominant gain occurs at $N\!=\!1{\to}2$ ($+92.4$ reward), after which performance saturates ($+0.6$ from $N\!=\!5$ to $N\!=\!20$): even a second client's experience is sufficient to cover CartPole's low-dimensional state space, pushing the policy close to the solved threshold.
On LunarLander, the critical transition occurs later at $N\!=\!2{\to}5$ ($+42.0$ reward)---its higher-dimensional, contact-rich dynamics require broader state-space coverage before federation yields a qualitative benefit.
Notably, LunarLander \emph{crosses} the solved threshold only at $N\!\geq\!5$ (rewards 223 and 232 vs.\ threshold 200), a discrete phase transition absent in the single- and two-agent settings (rewards 179 and 181).
CartPole approaches but does not consistently exceed its threshold at any tested $N$ within 600 episodes, suggesting that the solved criterion is more sensitive to per-episode sample efficiency than to federation breadth.

\textbf{Answer to Q3.} FedQHD scales gracefully with $N$, achieving up to $30\%$ improvement over independent learning with no degradation at large $N$.
Gains follow a diminishing-returns regime: a modest federation of $N\!=\!5$ agents already recovers most of the benefit and suffices to cross the solved threshold on LunarLander, consistent with the projection-residual bound in Proposition~1 saturating once the anchor feature matrix $X_i$ spans the dominant subspace of $Q^{\mathrm{glob}}_{\mathrm{ref}}$.

% \begin{figure}[t]
%     \centering
%     \includegraphics[width=0.8\linewidth]{figures/scalability.pdf}
%     \caption{Scalability: performance improvement vs. number of clients $N$.}
%     \label{fig:scalability}
% \end{figure}


\subsubsection{Ablation Study (Q4)}
\label{subsec:ablation}

We conduct two controlled ablations to validate the theoretical properties of the anchor-based aggregation step in FedQHD.
Both ablations use CartPole-v1 and measure the oracle Q-approximation error (Eq.~\ref{eq:metric_qaerr}): the RMSE between the ridge-regression estimate $\hat{Q}$ and the ground-truth $Q^*$ obtained from a converged reference agent ($D_\mathrm{ref}=2048$, $V(\pi_\mathrm{ref})=500$, 30 MC rollouts per anchor state).
Using $Q^*$ directly isolates approximation capacity from RL training noise, allowing clean verification of the theoretical scaling laws.
Each experiment is repeated over 3 seeds; each figure shows three panels: (left) smoothed RL learning curves (window 20), (centre) analytical Q-error vs.\ the varied parameter, and (right) final policy value gap $V(\pi) - V(\pi_\mathrm{ref})$.

% -------------------------------------------------
\paragraph{Ablation 1: Effect of Encoder Dimension $D$.}

\textbf{Setup.}
We vary $D \in \{16, 32, 64, 128, 256, 512, 1024, 2048\}$ while fixing $m = 4D$ (always in the over-determined regime, $m > D$), so the observed error reflects only the expressiveness of the RFF feature map rather than insufficient anchor coverage.
The RFF approximation theory~\cite{rahimi2007random} predicts the geometry floor scales as $\mathcal{O}(D^{-1/2})$.

\textbf{Results.}
Figure~\ref{fig:ablation1} confirms the predicted power-law decay.
The log-log slope is $\mathbf{-0.525}$, within $5\%$ of the theoretical $-0.5$, validating that FedQHD's anchor-based projection inherits the standard RFF approximation guarantee.
Quantitatively, RMSE decreases from $149.8 \pm 76.2$ at $D=16$ to $12.4 \pm 0.03$ at $D=2048$ — a $12\times$ reduction — with the steepest drop occurring between $D=16$ and $D=128$ (from $149.8$ to $26.4$, a $5.7\times$ gain) and diminishing returns beyond $D=512$ (from $14.6$ to $12.4$, only $1.2\times$).
The high variance at $D=16$ ($\pm 76.2$) reflects the instability of a very low-dimensional feature map: some random projections happen to align well with the CartPole Q-function while others do not.

The right panel reveals an important link between approximation quality and RL performance.
For $D \leq 64$, the shallow feature map yields RMSE $> 60$ and the RL policy fails to converge reliably ($V < 300$, high variance).
Once $D \geq 128$, RMSE drops below 27 and the policy consistently reaches $V \approx 500$ (the CartPole ceiling), indicating a practical sufficiency threshold: beyond $D=128$, additional dimensions reduce the theoretical geometry floor but no longer block policy convergence on this task.
This two-phase behaviour — sharp improvement then saturation — suggests that practitioners can use $D=512$--$1024$ as a cost-effective operating point, capturing over $97\%$ of the asymptotic accuracy at a fraction of the $D=2048$ cost.

\begin{figure}[t]
  \centering
  \includegraphics[width=\linewidth]{../results/figures/ablation1_vary_D.pdf}
  \caption{Ablation 1 — Effect of encoder dimension $D$ ($m=4D$ fixed).
           \emph{Left}: smoothed RL learning curves; higher $D$ converges faster and more reliably.
           \emph{Centre}: oracle Q-error on held-out anchors; log-log slope $-0.525$ (solid blue)
           closely matches the $D^{-1/2}$ theoretical bound (dashed red).
           \emph{Right}: final policy value gap; reliable convergence ($V\approx500$) requires $D\geq128$.}
  \label{fig:ablation1}
\end{figure}

% -------------------------------------------------
\paragraph{Ablation 2: Effect of Anchor Set Size $m$.}

\textbf{Setup.}
We fix $D=512$ and vary the number of anchor states $m \in \{51, 102, 256, 512, 1024, 2048\}$, corresponding to $m/D \in \{0.10, 0.20, 0.50, 1.00, 2.00, 4.00\}$.
Theory predicts a sharp regime transition at $m = D$:
\[
\mathrm{err}(m) \propto
\begin{cases}
\sqrt{D/m}, & m < D \;\text{(under-determined: system is rank-deficient)}, \\
D^{-1/2},   & m \gg D \;\text{(over-determined: geometry floor dominates)}.
\end{cases}
\]
In the under-determined regime the feature matrix $X \in \mathbb{R}^{m \times D}$ has fewer rows than columns, so many weight solutions $\hat{W}$ achieve zero training error while generalising poorly.
Once $m \geq D$, the normal equations become well-conditioned and the residual is governed solely by the RFF approximation quality of the encoder.

\textbf{Results.}
Figure~\ref{fig:ablation2} demonstrates the transition sharply and quantitatively.
In the under-determined regime ($m < D$), RMSE drops steeply with each doubling of $m$: from $5767 \pm 2104$ at $m/D=0.10$ to $618 \pm 152$ at $m/D=0.20$ ($9.3\times$) and to $91.2 \pm 15.9$ at $m/D=0.50$ ($6.8\times$).
Crossing the boundary to $m/D=1.00$ yields another $2.6\times$ reduction (to $35.1 \pm 1.0$), and beyond the transition the error stabilises near the geometry floor: $18.5 \pm 0.7$ at $m/D=2.00$ and $14.6 \pm 0.3$ at $m/D=4.00$.
The mean RMSE across under-determined settings ($m < D$) is $\mathbf{95\times}$ higher than across over-determined settings ($m \geq D$), confirming that $m = D$ is the critical threshold practitioners must exceed.

The right panel links the Q-error regimes directly to RL behaviour.
At $m/D=0.10$, the severely under-determined projection corrupts the aggregated weights sufficiently to destabilise learning: final $V = 353 \pm 177$, with large inter-seed variance indicating that some seeds fail entirely.
Even a modest increase to $m/D=0.20$ ($m=102$ anchors) eliminates this instability: $V = 500 \pm 0.0$ across all seeds, despite the Q-error still being $618$ — 42× above the geometry floor.
This reveals that RL policy quality is tolerant of absolute Q-error magnitude once the anchor system is not rank-deficient; what matters is whether the anchor set is sufficient to resolve the dominant action-value directions.
At $m/D=1.00$ ($m=D=512$) the policy shows renewed variance ($V=428 \pm 102$), likely because this is the boundary where the system transitions from one regime to the other and the ridge-regression solution is most sensitive to the specific anchor states sampled — a transient effect that disappears once $m > D$.

\textbf{Practical implication.}
Both ablations point to the same design rule: set $m \geq D$ and $D \geq 128$--$512$.
Under these conditions the anchor-based projection step introduces at most a $12\times$ RMSE overhead relative to the geometry floor, and RL convergence is robust across seeds.
For the main Q2 experiments we use $m=200$, $D_i \leq 10{,}000$; since $m < D_i$ for large clients, the under-determined regime analysis above explains why LunarLander Q2 shows greater performance degradation than CartPole Q2 (§\ref{subsec:performance_homo}).

\begin{figure}[t]
  \centering
  \includegraphics[width=\linewidth]{../results/figures/ablation2_vary_m.pdf}
  \caption{Ablation 2 — Effect of anchor set size $m$ ($D=512$ fixed).
           \emph{Left}: smoothed RL learning curves; all settings except $m/D=0.10$ converge to $V\approx500$.
           \emph{Centre}: oracle Q-error vs.\ $m/D$; the vertical dashed line marks the $m=D$ transition
           where the under-determined and over-determined regimes meet ($95\times$ gap in mean RMSE).
           \emph{Right}: final policy value; instability only at $m/D=0.10$ (rank-deficient regime).}
  \label{fig:ablation2}
\end{figure}

\textbf{Answer to Q4.}
Encoder dimension $D$ governs a hard geometry floor: Q-error decreases as $D^{-1/2}$, and $D \geq 128$ suffices for reliable RL convergence on the tested environments.
Anchor set size $m$ determines the projection regime: $m < D$ causes a catastrophic $95\times$ increase in approximation error relative to $m \geq D$.
Together, these results provide concrete design guidelines — $m \geq D \geq 512$ — that practitioners can use to configure FedQHD's heterogeneous aggregation step without empirical tuning.




%% Results
\section{Results}
\label{sec:results}

% Present experimental results
[Present the main experimental results with figures and tables.]

\subsection{Performance Comparison}
\label{subsec:performance}

[Compare FedQHD against baseline methods in terms of learning performance.]

% Example figure (uncomment and modify as needed):
% \begin{figure}[htbp]
%     \centering
%     \includegraphics[width=0.8\linewidth]{figures/performance_comparison.pdf}
%     \caption{Learning curves comparing FedQHD with baseline methods.}
%     \label{fig:performance}
% \end{figure}

\subsection{Communication Efficiency}
\label{subsec:communication}

[Analyze communication efficiency and overhead.]

\subsection{Scalability Analysis}
\label{subsec:scalability}

[Examine how the method scales with the number of clients or other parameters.]

\subsection{Ablation Study}
\label{subsec:ablation}

[Conduct ablation studies to understand the contribution of different components.]

% Example table (uncomment and modify as needed):
% \begin{table}[htbp]
%     \caption{Ablation study results}
%     \label{tab:ablation}
%     \centering
%     \begin{tabular}{lcc}
%         \hline
%         \textbf{Method} & \textbf{Metric 1} & \textbf{Metric 2} \\
%         \hline
%         Full Method & 0.85 & 0.90 \\
%         Without Component A & 0.78 & 0.85 \\
%         Without Component B & 0.80 & 0.88 \\
%         \hline
%     \end{tabular}
% \end{table}


%% Discussion
\section{Discussion}
\label{sec:discussion}

% Interpret and discuss the results
[Interpret the experimental results and discuss their implications.]

\subsection{Key Findings}
\label{subsec:key_findings}

[Summarize the key findings from the experiments.]

\subsection{Limitations}
\label{subsec:limitations}

[Discuss limitations of the proposed approach and experiments.]

\subsection{Future Directions}
\label{subsec:future_work}

[Suggest potential future research directions.]


%% Conclusion
\section{Conclusion}
\label{sec:conclusion}

% Summarize the paper
[Summarize the main contributions and findings of the paper.]

[Restate the importance of the problem and how your approach addresses it.]

[Highlight the key experimental results that support your claims.]

[Conclude with the broader impact or significance of this work.]


%%%%%%%%%%%%%%%%%%%%%%%%%%%%%%%%%%%%%%%%%%%%%%%%%%%%%%%%%%%%%%%%
%% FORMATTING REMINDER: Keep the following guidelines in mind
%%%%%%%%%%%%%%%%%%%%%%%%%%%%%%%%%%%%%%%%%%%%%%%%%%%%%%%%%%%%%%%%
% Citations: Use \citet{} for in-text citations and \citep{} for parenthetical
% Figures: Always reference figures before they appear; use \ref{fig:label}
% Tables: Caption goes BEFORE the table; use \ref{tab:label}
% Equations: Use align environment with & after = or inequality symbols
% Sections: Use \section, \subsection, and \subsubsection for structure

%%%%%%%%%%%%%%%%%%%%%%%%%%%%%%%%%%%%%%%%%%%%%%%%%%%%%%%%%%%%%%%%
%% Optional: Broader Impact Statement
%%%%%%%%%%%%%%%%%%%%%%%%%%%%%%%%%%%%%%%%%%%%%%%%%%%%%%%%%%%%%%%%
% \section*{Broader Impact Statement}
% \label{sec:broaderImpact}
% Discuss possible repercussions of your work and any potential negative impact. 

%%%%%%%%%%%%%%%%%%%%%%%%%%%%%%%%%%%%%%%%%%%%%%%%%%%%%%%%%%%%%%%%
%% Appendices (Optional)
%%%%%%%%%%%%%%%%%%%%%%%%%%%%%%%%%%%%%%%%%%%%%%%%%%%%%%%%%%%%%%%%
% \appendix
%
% \section{Additional Experimental Details}
% \label{sec:appendix_experiments}
% [Include additional experimental details, hyperparameters, etc.]
%
% \section{Proofs}
% \label{sec:appendix_proofs}
% [Include theoretical proofs if applicable]
%
% \subsubsection*{Acknowledgments}
% \label{sec:ack}
% [Add acknowledgments here after paper acceptance and deanonymization]

%%%%%%%%%%%%%%%%%%%%%%%%%%%%%%%%%%%%%%%%%%%%%%%%%%%%%%%%%%%%%%%%
%% NOTE: THIS MARKS THE END OF THE "MAIN TEXT"
%%%%%%%%%%%%%%%%%%%%%%%%%%%%%%%%%%%%%%%%%%%%%%%%%%%%%%%%%%%%%%%%

%%%%%%%%%%%%%%%%%%%%%%%%%%%%%%%%%%%%%%%%%%%%%%%%%%%%%%%%%%%%%%%%
%% Bibliography
%%%%%%%%%%%%%%%%%%%%%%%%%%%%%%%%%%%%%%%%%%%%%%%%%%%%%%%%%%%%%%%%
\bibliography{main}
\bibliographystyle{rlj}

%%%%%%%%%%%%%%%%%%%%%%%%%%%%%%%%%%%%%%%%%%%%%%%%%%%%%%%%%%%%%%%%
% AUTHOR: If your paper has no supplementary materials, you may 
%         comment out the line below, which creates the title for
%         the supplementary materials.
%%%%%%%%%%%%%%%%%%%%%%%%%%%%%%%%%%%%%%%%%%%%%%%%%%%%%%%%%%%%%%%%
\beginSupplementaryMaterials

Content that appears after the references are not part of the ``main text,'' have no page limits, are not necessarily reviewed, and should not contain any claims or material central to the paper. 
%
If your paper includes supplementary materials, use the \begin{center}
    {\tt {\textbackslash}beginSupplementaryMaterials} 
\end{center}
command as in this example, which produces the title and disclaimer above. 
%
If your paper does not include supplementary materials, this command can be removed or commented out.

\end{document}
